On lance une pierre, de masse $m=\qty{200}{\g}$, verticalement et vers le haut
avec une vitesse de valeur $v=\qty{5.0}{\m\per\s}$.
\begin{enumerate}
    \item Calculer l'énergie cinétique et l'énergie potentielle de pesanteur de
          la pierre au moment du départ (on considérera qu'elle est initialement
          au niveau du sol).
    \item On néglige les frottements de l'air sur la pierre. Quelle
          transformation d'énergie se produit lors de l'ascension de la pierre~?
    \item Caractériser la position pour laquelle l'énergie potentielle de
          pesanteur de la pierre est maximale. Quelle est alors sa valeur~?
    \item Jusqu'à quelle hauteur la pierre s'élève-t-elle au-dessus du sol~?
    \item Quelle est la vitesse de la pierre lorsqu'elle revient au sol~?
\end{enumerate}

